\section{Technical Supplement: Mathematical Foundations of OPH}\label{technical-supplement-mathematical-foundations-of-observer-patch-holography}

\begin{quote}
OPH is an observer-centric reconstruction program for fundamental physics. In the current implementation it uses two external inputs (pixel area and screen capacity) together with structural axioms and explicit regulator/scaling premises. This supplement collects the mathematical ingredients for the conditional Lorentz/null-modular/Einstein branch, conditional compact gauge reconstruction, and the related bosonic quantitative program.
\end{quote}

\subsection{For Theoretical Physicists}\label{for-theoretical-physicists}

This document collects mathematical derivations and stated premises supporting the OPH program. The Lorentz/null-modular/Einstein and compact-gauge statements are scaling-limit / categorical results under explicit regularity hypotheses; these hypotheses are target properties of the intended EFT phase and candidate UV completions, not fixed-cutoff matrix-algebra theorems.

\textbf{Scope convention.} Statements about Lorentz recovery, modular flow, and the Einstein branch are scaling-limit / EFT statements rather than literal fixed-cutoff type-I claims. Statements about gauge reconstruction refer to the refinement-stable directed colimit of fixed-cutoff sector categories, not to the finite block list at one regulator scale.

\begin{center}\rule{0.5\linewidth}{0.5pt}\end{center}

\section{1. Core Axioms and Definitions}\label{1-core-axioms-and-definitions}

\subsection{1.1 Canonical Five-Axiom Package}\label{11-the-axiom-set}

This supplement uses the same authoritative OPH premise ledger as the compact note and the main manuscript. The canonical axioms are:

\textbf{Screen Net:} Physical reality is encoded on a horizon screen \(S^2\) carrying a net of von Neumann algebras \(P \mapsto \mathcal{A}(P)\) for patches \(P \subset S^2\).

\textbf{Overlap Consistency:} For overlapping patches \(P_1 \cap P_2 \neq \emptyset\), the restrictions of local states must agree: \[\omega_1|_{\mathcal{A}(P_1 \cap P_2)} = \omega_2|_{\mathcal{A}(P_1 \cap P_2)}\]

\textbf{Local MaxEnt and Refinement Stability:} At the UV scale the realized state maximizes entropy subject to a finite family of gauge-invariant local constraints, and the resulting family is stable under coarse-graining.

\textbf{Recoverable Generalized Entropy:} A generalized entropy functional exists: \[S_{\text{gen}}(C) = \frac{A(\partial C)}{4G} + S_{\text{bulk}}(C)\] satisfying quantum focusing and the recoverability structure used in collar and null-strip arguments.

\textbf{Minimal Admissible Realization (MAR):} On the admissible low-energy branch, the realized sector package is the lexicographically minimal one under the complexity vector \(C(\mathfrak S)\). MAR enters the gauge-selection branch; the Lorentz/null-modular and Einstein discussions below mainly use the first four axioms plus the technical premises stated where needed.

\subsection{1.2 Technical Premises and Legacy Regulator Names}\label{12-regulator-premises}

The supplement continues to use historical labels such as R0, R1, LR, I, and N1--N3 because later derivations cite them directly. In the present version, R0/R1 are not extra microphysical axioms, and LR together with the refinement-stable qualifier used later are not imported regularity language. On the local finite-constraint MaxEnt branch, they are the propagation and persistence consequences of the canonical third axiom.

\textbf{R0 (finite type-I regulator presentation):} On a fixed UV cellulation, each cell carries finite local data and can be Hilbertized on a finite-dimensional space \(\tilde{\mathcal H}_i\). For any finite union \(R\) of cells, the extended algebra before quotienting by overlap redundancy is
\[
\widetilde{\mathcal A}(R)=\mathcal B(\tilde{\mathcal H}_R),
\qquad
\tilde{\mathcal H}_R=\bigotimes_{i\subset R}\tilde{\mathcal H}_i.
\]

\textbf{R1 (boundary fixed-point realization of overlap gluing):} The redundancy in overlap identifications acts only on cut data. On each finite-dimensional chart, its unitary image has compact closure, so for a region \(R \subset S^2\) one may represent it by a compact boundary group \(G_{\partial R}\) acting on \(\tilde{\mathcal H}_R\), with physical algebra
\[
\mathcal{A}(R)=\mathcal{B}(\tilde{\mathcal H}_R)^{G_{\partial R}}.
\]

\textbf{Derived local-Gibbs branch.} If the UV constraints are expectations of finitely many gauge-invariant local operators \(\{O_a(x)\}\) of support radius \(O(\ell_{\mathrm{UV}})\), then the entropy maximizer is
\[
\omega_\ell(\lambda)= Z_\ell(\lambda)^{-1}\exp\!\left(-K_\ell(\lambda)\right),
\qquad
K_\ell(\lambda)=\sum_x \sum_a \lambda_a O_a(x)+\sum_b \mu_b Q_b.
\]
Thus the logarithm of the state is a quasi-local UV generator built from the same bounded-support densities that define the branch.

\textbf{Derived LR / T2 bound.} For local observables \(A_X,B_Y\) and the automorphism group \(\tau_t^\ell\) generated by \(K_\ell(\lambda)\), or by any branch generator in the same bounded-support algebraic closure, standard Lieb--Robinson estimates give
\[
\bigl\|[\tau_t^\ell(A_X),B_Y]\bigr\|
\le
C\,\|A_X\|\,\|B_Y\|\,\min(|X|,|Y|)\,
e^{-(d(X,Y)-v_\ell|t|)/\xi}.
\]
So finite-velocity quasi-local propagation is internal to the local MaxEnt branch rather than an extra premise.

\textbf{Derived endpoint control / T3.} After removal of the central endpoint term, changing a bounded interval \(I\) to \(I'\) modifies only an \(O(|I'\Delta I|)\) collar of local terms. Hence on any fixed local-energy-bounded domain
\[
\left|\langle\psi,(K[I']-K[I])\phi\rangle\right|
\le
C_{\psi,\phi,I_{\max}}\,|I'\Delta I|.
\]
This is the bounded-interval endpoint-Lipschitz / finite-variation control later used in the null modular derivation.

\textbf{Derived refinement-stable branch.} Because the same finite constraint family is preserved under coarse-graining, the regulator states belong to one common finite-dimensional multiplier family. The realized branch is the stable or fixed branch inside that family. Whenever later sections speak of a refinement-stable directed colimit of edge sectors, the refinement-stable qualifier means persistence along this branch; only transportability, symmetry, bosonic, and fiber-functor clauses remain extra when gauge reconstruction is invoked.

\textbf{Read-along convention.} Older references to ``A3'' mean Recoverable Generalized Entropy; older references to ``A4'' mean the recoverability clause used together with the explicit theorem-local collar or null-strip hypotheses stated below, not a separate extra-premise package; older references to ``LR'' mean the derived quasi-local propagation bound above; older references to ``Assumption I'' mean the refinement-stable branch condition already contained in the canonical third axiom. The remaining nontrivial burden is therefore the controlled refinement limit and the later geometric or null modular premises, not the matrix or fixed-point bookkeeping, quasi-local propagation, or branch stability at one regulator scale.

\subsection{1.3 Edge-Center Completion (EC)}\label{13-edge-center-completion-ec}

Older supplement drafts treated the boundary gluing group as part of the regulator package and then concluded that EC implies exact Markovity. The sharper statement used in the compact note is weaker and more precise. At fixed regulator scale, overlap consistency itself determines the boundary redundancy. On the ordinary or central-defect branch, EC is the invariant-theoretic consequence of that derived boundary action. Exact Markovity is a further state property and is not implied by EC alone. This is the reading rule for the historical \(R0/R1\) shorthand.

\textbf{Proposition 1.1a (Derived boundary gluing datum).} Choose a finite regulator chart for the patches meeting along a connected cut \(\Sigma\). Because the local overlap algebras are finite-dimensional matrix algebras, any overlap-consistent recharting is an inner automorphism and is implemented by a unitary on the cut Hilbert space. The compact closure of the subgroup generated by these recharting unitaries is a compact boundary redundancy group \(K_\Sigma\). If triple-overlap defects are central, the projective composition law lifts to a compact central extension \(\widehat K_\Sigma\); on the ordinary branch one simply sets \(\widehat K_\Sigma = K_\Sigma\). A genuinely noncentral 2-group defect is the only obstruction to reducing the overlap transition system to an ordinary compact group action.

\textbf{Proof.} On a finite-dimensional matrix algebra every \(^*\)-automorphism is inner, so each overlap-consistent recharting is implemented by a unitary on the cut Hilbert space. The subgroup generated by those unitaries has compact closure inside a finite-dimensional unitary group. A central defect gives a projective composition law that lifts to a central extension; a genuinely noncentral defect is the remaining fixed-cutoff obstruction. QED.

\textbf{Theorem 1.1 (Derived EC decomposition).} Under the fixed-cutoff regulator realization of Section 1.2, and on the ordinary or central-defect branch, the collar Hilbert space is
\[
\mathcal H_{B_\delta}
=
(\tilde{\mathcal H}_{B_L}\otimes \tilde{\mathcal H}_{B_R})^{\widehat K_\Sigma}
\cong
\bigoplus_{\alpha}
\left(\mathcal H_{b_L^\alpha}\otimes \mathcal H_{b_R^\alpha}\right),
\]
and the center of the collar algebra is generated by the block projectors:
\[
Z(\mathcal A(B_\delta))
=
\bigoplus_\alpha \mathbb C\cdot \mathbf 1_\alpha.
\]
The right half-collar carries the contragredient representation because it sees the inverse transport across the same cut.

\textbf{Proof.} Decompose the left boundary data into irreps
\[
\tilde{\mathcal H}_{B_L}
=
\bigoplus_\alpha
\left(W_\alpha\otimes \mathcal H_{b_L^\alpha}\right),
\]
and the right boundary data into dual irreps
\[
\tilde{\mathcal H}_{B_R}
=
\bigoplus_\beta
\left(W_\beta^*\otimes \mathcal H_{b_R^\beta}\right).
\]
Then
\[
\tilde{\mathcal H}_{B_L}\otimes \tilde{\mathcal H}_{B_R}
=
\bigoplus_{\alpha,\beta}
(W_\alpha\otimes W_\beta^*)\otimes
(\mathcal H_{b_L^\alpha}\otimes \mathcal H_{b_R^\beta}).
\]
By Schur's lemma,
\[
(W_\alpha\otimes W_\beta^*)^{\widehat K_\Sigma}
\cong
\begin{cases}
\mathbb C, & \alpha=\beta,\\
0, & \alpha\neq \beta,
\end{cases}
\]
so only matched left/right labels survive, and the surviving block projectors generate the center. QED.

\textbf{Corollary 1.2 (Exact Markov adds the state factorization).} If, in addition,
\[
I_\omega(A_\delta:D_\delta\mid B_\delta)=0,
\]
or one passes to the explicitly stated idealized recoverability limit that reduces to exact Markovity, then
\[
\rho_{A_\delta B_\delta D_\delta}
=
\bigoplus_\alpha p_\alpha
\left(\rho_{A_\delta b_L^\alpha}\otimes \rho_{b_R^\alpha D_\delta}\right).
\]
EC therefore gives the kinematic block decomposition; exact Markovity is the extra state input that gives the HJPW normal form.

\textbf{Remark.} Approximate recoverability gives controlled deviations from this normal form; it is not implied by EC alone. The only remaining fixed-cutoff blocker to an ordinary-group EC theorem is a genuinely noncentral 2-group defect. The later open burden is the controlled refinement or scaling regime and the modular or null-limit premises, not the finite-dimensional collar decomposition itself.
\subsection{1.4 Current Continuous Inputs}\label{14-fundamental-parameters}

In the current quantitative implementation, OPH uses exactly two external continuous inputs:

\begin{enumerate}
\def\labelenumi{\arabic{enumi}.}
\tightlist
\item
  \textbf{Pixel area:} \(a_{\text{cell}} \approx 1.63 \ell_P^2\)
\item
  \textbf{Screen capacity:} \(\log(\dim \mathcal{H}_{\text{tot}}) \sim 10^{122}\)
\end{enumerate}

From these:

\begin{itemize}
\tightlist
\item
  Newton\textquotesingle s constant: \(G = \frac{a_{\text{cell}}}{4\bar{\ell}(t)}\)
\item
  Cosmological constant: \(\Lambda = \frac{3\pi}{G \cdot \log(\dim \mathcal{H}_{\text{tot}})}\)
\end{itemize}

\begin{center}\rule{0.5\linewidth}{0.5pt}\end{center}

\section{2. Emergence of Gravity and Quantum Mechanics}\label{2-emergence-of-gravity-and-quantum-mechanics}

\subsection{2.1 Local Lorentz Structure}\label{21-local-lorentz-structure}

\textbf{Theorem 2.1 (Conformal Group Isomorphism):} The orientation-preserving conformal group of \(S^2\) is isomorphic to the connected Lorentz group: \[\text{Conf}^+(S^2) \cong \text{PSL}(2,\mathbb{C}) \cong \text{SO}^+(3,1)\]

This provides the kinematic structure of special relativity.

\subsection{2.2 Geometric Modular Flow}\label{22-geometric-modular-flow}

The collar technology determines the controlled approach to an exact Markov normal form, but by itself it does not force a specific continuum modular geometry. The theorem-scope must therefore be stated as an OPH branch statement. In the attached source set, the sharp claim is: if the refinement-stable MaxEnt branch realized by OPH lies in the Bisognano--Wichmann geometric modular class for caps, then that branch already supplies the standard cap-preserving conformal dilation and the theorem fixes only its \(2\pi\) normalization plus the carried collar remainder. This replaces the earlier packaging that treated cap isotropy and a separate Euclidean-period postulate as independent premises.

For a shrinking collar family \((A_\delta,B_\delta,D_\delta)\) around \(\partial C\), write
\[
\varepsilon_\delta:=I(A_\delta:D_\delta\mid B_\delta)_\omega,
\qquad
r_{\mathrm{FR}}(\varepsilon_\delta):=
2\sqrt{1-e^{-\varepsilon_\delta}}
\le 2\sqrt{\varepsilon_\delta},
\]
and, when the relevant marginals are uniformly faithful with lower spectral bound \(\lambda_\ast>0\),
\[
\eta_\delta^{\mathrm M}
:=
4\lambda_\ast^{-1}\,
\delta^{\mathrm M}_{A_\delta:B_\delta:D_\delta}(\varepsilon_\delta).
\]
All exact collar formulas in this subsection are therefore to be read either literally at exact Markovity or along a controlled collar family satisfying
\[
\delta^{\mathrm M}_{A_\delta:B_\delta:D_\delta}(\varepsilon_\delta)\to 0
\qquad(\delta\downarrow 0),
\]
with \(r_{\mathrm{FR}}(\varepsilon_\delta)\) and \(\eta_\delta^{\mathrm M}\) carried explicitly.

\textbf{Theorem 2.2 (Geometric modular flow on caps in the OPH Bisognano--Wichmann branch).} Assume the OPH axioms, the regulator premises R0--R1, local MaxEnt/refinement stability, the collar-control hypotheses used earlier in the supplement, and the branch condition that the refinement-stable cap net belongs to the Bisognano--Wichmann geometric modular class for caps and their conformal images, so that for each cap \(C\subset S^2\) the limiting modular flow is implemented by the standard cap-preserving conformal one-parameter subgroup \(\lambda_C(s)\), with geometric parameter \(s\) normalized by the null blow-up \(v\mapsto e^{-s}v\). If \(B_C\) denotes its infinitesimal generator, then
\[
K_C=2\pi B_C.
\]

More precisely, at collar scale \(\delta\),
\[
K_C^{(\delta)}=2\pi B_C+E_C^{(\delta)},
\]
where the carried remainder is controlled on bounded observables by \(r_{\mathrm{FR}}(\varepsilon_\delta)\) and on modular-Hamiltonian replacements by \(\eta_\delta^{\mathrm M}\). Hence \(E_C^{(\delta)}\to 0\) on every controlled collar family whose faithfulness bound survives the refinement limit.

\textbf{Proof.} The recoverability/collar analysis localizes the modular defect to the shrinking collar and allows replacement of the physical state by an exact Markov reference with explicit carried error. On the chosen BW refinement branch, the standard cap-preserving conformal subgroup \(\lambda_C(s)\) is part of the branch content, so the limiting modular Hamiltonian is \(K_C=\kappa B_C\) for some \(\kappa>0\). The modular flow is KMS by definition, and because the geometric parameter \(s\) is normalized by unit null dilation on the branch, the usual Bisognano--Wichmann relation fixes \(\kappa=2\pi\). The displayed remainder \(E_C^{(\delta)}\) is exactly the carried collar-replacement error, so it vanishes in the controlled refinement limit. QED.

\textbf{Scope note.} The fixed-cutoff OPH collar technology establishes the \(2\pi\) normalization and the carried collar remainder on the Bisognano--Wichmann geometric modular branch stated above. The branch-selection question is encoded in the stated hypotheses for that branch.

\subsection{2.3 Null Modular Bridge (EFT Conditions N1-N3)}\label{23-null-modular-bridge-eft-conditions-n1-n3}

The fixed-cutoff null bridge reaches the derived half-sided-inclusion step. What is established here is: transferred cut-center data on null strips, exact-or-controlled strip additivity on one inherited strip model, endpoint-Lipschitz control of the renormalized half-line family, and the geometric blow-up argument that turns the null half-line net into a half-sided modular pair, so that Borchers--Wiesbrock yields an explicit positive null-translation generator on its Stone domain. What is not established here is the downstream density upgrade or the relativistic null-stress identification. The present condition package should therefore be read with the following split.

\textbf{Condition N1 (null-strip center transfer and exact-or-controlled strip additivity).} For consecutive null strips
\[
I_-=(v_1,v_2),\qquad
J=(v_2,v_3),\qquad
I_+=(v_3,v_4),
\]
assume the two cuts inherit the same compact overlap-redundancy data as ordinary collars, so that on a compatible type-I regulator presentation
\[
\tilde{\mathcal H}_{J}
\cong
\bigoplus_{\alpha_-,\alpha_+}
W_{\alpha_-}^{*}\otimes
\mathcal H_{j^{\alpha_-,\alpha_+}}\otimes
W_{\alpha_+}
\]
and the strip algebra is the commutant of the transported cut action. Then
\[
\mathcal A(J)
\cong
\bigoplus_{\alpha_-,\alpha_+}
\mathcal B(\mathcal H_{j^{\alpha_-,\alpha_+}}),
\qquad
Z(\mathcal A(J))
=
\bigoplus_{\alpha_-,\alpha_+}\mathbb C\,P_{\alpha_-,\alpha_+}.
\]
If, in addition, each multiplicity space admits an inherited split
\[
\mathcal H_{j^{\alpha_-,\alpha_+}}
\cong
\mathcal H_{j_L^{\alpha_-,\alpha_+}}
\otimes
\mathcal H_{j_R^{\alpha_-,\alpha_+}},
\]
with \(\mathcal A(I_-\cup J)\) acting blockwise only on
\(\mathcal H_{i_-^{\alpha_-}}\otimes \mathcal H_{j_L^{\alpha_-,\alpha_+}}\)
and \(\mathcal A(J\cup I_+)\) acting blockwise only on
\(\mathcal H_{j_R^{\alpha_-,\alpha_+}}\otimes \mathcal H_{i_+^{\alpha_+}}\),
then exact Markov strip states satisfy exact four-term modular additivity on the canonical inherited model:
\[
\Delta K_J(\omega)
:=
K_{I_-\cup J\cup I_+}(\omega)-K_{I_-\cup J}(\omega)-K_{J\cup I_+}(\omega)+K_J(\omega)
=0,
\]
equivalently up to a central endpoint term. For a general strip state with
\[
I(I_-:I_+\mid J)_\omega\le \varepsilon,
\]
define
\[
\mathfrak M_{I_-:J:I_+}
:=
\left\{
\sigma:\ I(I_-:I_+\mid J)_\sigma=0
\right\},
\qquad
\delta^{\mathrm M}_{I_-:J:I_+}(\varepsilon)
:=
\sup\left\{
\inf_{\sigma\in\mathfrak M_{I_-:J:I_+}}\|\rho-\sigma\|_1:
I(I_-:I_+\mid J)_\rho\le\varepsilon
\right\}.
\]
Choose \(\widetilde\omega_J\in \mathfrak M_{I_-:J:I_+}\) with
\[
\|\omega-\widetilde\omega_J\|_1
\le
\delta^{\mathrm M}_{I_-:J:I_+}(\varepsilon),
\]
and set
\[
\mathfrak D_J(\omega,\widetilde\omega_J)
:=
\Delta K_J(\omega)-\Delta K_J(\widetilde\omega_J).
\]
Then
\[
K_{I_-\cup J\cup I_+}(\omega)
=
K_{I_-\cup J}(\omega)+K_{J\cup I_+}(\omega)-K_J(\omega)
+K_{\partial,J}(\widetilde\omega_J)
+\mathfrak D_J(\omega,\widetilde\omega_J),
\]
with \(K_{\partial,J}(\widetilde\omega_J)\) central. If the relevant marginals are uniformly faithful with lower spectral bound \(\lambda_\ast>0\), then
\[
\|\mathfrak D_J(\omega,\widetilde\omega_J)\|_\infty
\le
4\lambda_\ast^{-1}\,
\delta^{\mathrm M}_{I_-:J:I_+}(\varepsilon).
\]
Independently, Fawzi--Renner recovery supplies a constructive comparison state with
\[
r_{\mathrm{FR}}(\varepsilon)
=
2\sqrt{1-e^{-\varepsilon}}
\le
2\sqrt{\varepsilon}.
\]
Thus N1 is exact only at exact Markovity, or else along a controlled strip family on one fixed inherited strip model with
\[
\delta^{\mathrm M}_{I_-:J:I_+}(\varepsilon_J)\to 0.
\]

\textbf{Condition N2 (derived half-sided modular inclusion on the null half-line blow-up net).} Fix a smooth cut point, choose affine coordinate \(v\) on the chosen null generator \(\Omega\) so that the cut sits at \(v=0\), and define for \(a\ge 0\)
\[
H_a:=(a,\infty),
\qquad
\mathcal M_a(\Omega)
:=
\overline{\bigvee_{a<c<d<\infty}\mathcal A((c,d),\Omega)}.
\]
By isotony of the null interval net,
\[
a\le b
\quad\Longrightarrow\quad
\mathcal M_b(\Omega)\subseteq \mathcal M_a(\Omega).
\]
On the OPH geometric scaling branch, the blow-up of cap modular flow acts by the null dilation
\[
v-v_0\mapsto e^{-2\pi t}(v-v_0).
\]
Hence
\[
\sigma_t^\omega\bigl(\mathcal M_a(\Omega)\bigr)
=
\mathcal M_{e^{-2\pi t}a}(\Omega).
\]
For \(t\le 0\) one has \(e^{-2\pi t}a\ge a\), so \(H_{e^{-2\pi t}a}\subseteq H_a\) and therefore
\[
\sigma_t^\omega\bigl(\mathcal M_a(\Omega)\bigr)\subseteq \mathcal M_a(\Omega).
\]
Thus \(\mathcal M_a(\Omega)\subset \mathcal M_0(\Omega)\) is half-sided modular. After the harmless convention change \(t\mapsto -t\), this is the standard positive-time form. The old N2 premise is therefore derived from OPH null-interval structure, isotony, and the scaling-limit geometric action rather than imported separately.

\textbf{Condition N3 (endpoint-Lipschitz control and weak tail generator).} For renormalized modular Hamiltonians
\[
\widetilde K[I,\Omega]
:=
K[I,\Omega]-K_\partial(I,\Omega),
\qquad
\widetilde K_a(\Omega):=\widetilde K[(a,\infty),\Omega],
\]
bounded intervals satisfy
\[
\bigl|
\langle\psi,(\widetilde K[(a',b'),\Omega]-\widetilde K[(a,b),\Omega])\phi\rangle
\bigr|
\le
C_{\psi,\phi,\Omega}\bigl(|a'-a|+|b'-b|\bigr),
\]
and half-lines satisfy
\[
\bigl|
\langle\psi,(\widetilde K_{a'}(\Omega)-\widetilde K_a(\Omega))\phi\rangle
\bigr|
\le
C_{\psi,\phi,\Omega}|a'-a|.
\]
Hence, for fixed \(\psi,\phi\),
\[
f_{\psi,\phi}(a):=\langle\psi,\widetilde K_a(\Omega)\phi\rangle
\]
is locally Lipschitz and absolutely continuous, and its distributional derivative defines the weak tail generator
\[
\langle\psi,q(a,\Omega)\phi\rangle
:=
-\frac{1}{2\pi}\,\partial_a f_{\psi,\phi}(a).
\]
For \(a<b\),
\[
\langle\psi,(\widetilde K_b(\Omega)-\widetilde K_a(\Omega))\phi\rangle
=
-2\pi\int_a^b \langle\psi,q(v,\Omega)\phi\rangle\,dv.
\]
This is the end of the fixed-cutoff analytic part of the bridge. The stronger assumptions
\[
q \ \text{weakly differentiable},
\qquad
f_{\psi,\phi}(a)\to 0,
\qquad
q_{\psi,\phi}(a)\to 0
\quad (a\to+\infty)
\]
would permit the downstream density formula
\[
\langle\psi,p(a,\Omega)\phi\rangle
:=
-\partial_a\langle\psi,q(a,\Omega)\phi\rangle
=
\frac{1}{2\pi}\,\partial_a^2 f_{\psi,\phi}(a),
\qquad
\langle\psi,\widetilde K_a(\Omega)\phi\rangle
=
2\pi\int_a^\infty (v-a)\,\langle\psi,p(v,\Omega)\phi\rangle\,dv,
\]
but that density upgrade is not established here.

\textbf{Imported Borchers consequence.} Let \(\Delta_0(\Omega)\) be the modular operator of the standard pair \((\mathcal M_0(\Omega),\omega)\), define \(K_0(\Omega):=-\log \Delta_0(\Omega)\), and let \(\Delta_a(\Omega)\) denote the modular operator of \((\mathcal M_a(\Omega),\omega)\). Applying Borchers--Wiesbrock to the derived half-sided pair of Condition N2 yields a unique strongly continuous unitary group \(U_\Omega(a)=e^{iaP_\Omega}\) with positive self-adjoint Stone generator \(P_\Omega\) on
\[
D(P_\Omega)
:=
\left\{
\psi\in\mathcal H:
\lim_{a\to0}\frac{U_\Omega(a)\psi-\psi}{ia}\ \text{exists}
\right\}.
\]
It obeys
\[
\Delta_0(\Omega)^{it}U_\Omega(a)\Delta_0(\Omega)^{-it}
=
U_\Omega(e^{-2\pi t}a),
\qquad
[K_0(\Omega),P_\Omega]=-\,i\,2\pi P_\Omega,
\]
and if \(K_a(\Omega):=-\log \Delta_a(\Omega)\), then
\[
K_a(\Omega)=K_0(\Omega)-2\pi a\,P_\Omega
\]
as a quadratic-form identity on \(D(K_0(\Omega))\cap D(P_\Omega)\).

\textbf{Consequence.} The present condition package proves the fixed-cutoff strip bridge through the weak tail generator \(q(a,\Omega)\), derives half-sided modular inclusion on the null half-line blow-up net, and therefore identifies an explicit positive null-translation generator \(P_\Omega\) together with its Stone domain and half-line modular-Hamiltonian relation. The remaining downstream burdens are the density-upgrade hypotheses and the relativistic EFT null-stress premise. Those later steps are explicit rather than hidden parts of N1--N3.
\subsection{2.4 Stress Tensor Reconstruction}\label{24-stress-tensor-reconstruction}

\textbf{Lemma 2.3 (Null Data Ambiguity):} If a symmetric tensor \(X_{ab}\) satisfies \(X_{ab}k^a k^b = 0\) for all null \(k\), then: \[X_{ab} = \phi \, g_{ab}\] for some scalar \(\phi\).

\emph{Proof:} Decompose \(X_{ab} = Y_{ab} + \phi g_{ab}\) with \(Y\) traceless. For null \(k = (1, \hat{n})\): \[X_{kk} = Y_{kk} + \phi g_{kk} = Y_{kk}\] since \(g_{kk} = 0\). Expanding: \[Y_{kk} = Y_{00} + 2\hat{n}^i Y_{0i} + \hat{n}^i \hat{n}^j Y_{ij}\]

If this vanishes for all \(\hat{n} \in S^2\), each angular moment (scalar, vector, traceless tensor) vanishes independently, so \(Y_{ab} = 0\). \(\square\)

\textbf{Corollary 2.4:} Null modular data determine \(T_{ab}\) only up to \(\phi g_{ab}\). Consequently, the Einstein equation is fixed only up to \(\Lambda g_{ab}\).

\textbf{Reconstruction Formulas:} From \(f(\hat{n}) := T_{kk}(\hat{n})\): \[T_{0i} = \frac{3}{2}\langle \hat{n}_i f \rangle_{S^2}\] \[T_{ij} - \frac{\delta_{ij}}{3}T^{(\text{spatial})}_{kk} = \frac{15}{2}\left\langle \left(\hat{n}_i \hat{n}_j - \frac{\delta_{ij}}{3}\right) f \right\rangle_{S^2}\]

\subsection{2.5 Entanglement Equilibrium and Einstein\textquotesingle s Equation}\label{25-entanglement-equilibrium-and-einsteins-equation}

\textbf{Theorem 2.5 (Conditional Jacobson-type Derivation):} In the same local Lorentzian scaling regime, under MaxEnt state selection with generalized entropy stationarity, the null-stress identification premise of the relativistic EFT branch, and admissible first variations about a maximally symmetric reference state,
\[\delta S_{\text{gen}} = 0\]

implies the rest-frame first-variation relation
\[
\delta\!\left(G_{00} + \Lambda g_{00}\right) = 8\pi G \,\delta\langle T_{00} \rangle.
\]

If this first-variation relation holds for all local directions and reference states in the scaling regime, overlap consistency upgrades it to the full tensor equation
\[
G_{ab} + \Lambda g_{ab} = 8\pi G \langle T_{ab} \rangle
\]
modulo the expected \(\Lambda g_{ab}\) ambiguity.

\textbf{Derivation sketch:}

\begin{enumerate}
\def\labelenumi{\arabic{enumi}.}
\item
  Raychaudhuri equation relates area change to \(R_{kk}\): \[\frac{d\theta}{d\lambda} = -\frac{1}{2}\theta^2 - \sigma^2 - R_{kk}\]
\item
  Recoverable Generalized Entropy implies: \[\frac{d}{d\lambda}\left(\frac{A}{4G} + S_{\text{bulk}}\right) \leq 0\]
\item
  Stationarity at MaxEnt gives: \[R_{kk} = 8\pi G \langle T_{kk} \rangle\]
\item
  By Lemma 2.3, this fixes Einstein\textquotesingle s equation up to \(\Lambda g_{ab}\).
\end{enumerate}

\subsection{2.6 Newton\textquotesingle s Constant from Edge Entropy}\label{26-newtons-constant-from-edge-entropy}

\textbf{Theorem 2.6:} The gravitational coupling is fixed by the edge entropy density: \[G = \frac{a_{\text{cell}}}{4\bar{\ell}(t)}\]

where \(\bar{\ell}(t) = a_{\text{cell}}/(4\ell_P^2)\) is the dimensionless entropy per cell.

With \(a_{\text{cell}} \approx 1.63 \ell_P^2\): \[\bar{\ell} = \frac{1.63}{4} \approx 0.4077\]

\begin{center}\rule{0.5\linewidth}{0.5pt}\end{center}

\section{3. The Measurement Problem}\label{3-the-measurement-problem}

\subsection{3.1 Records as Central Variables}\label{31-records-as-central-variables}

\textbf{Definition 3.1:} A \emph{record} in patch \(P\) is a family of orthogonal projectors \(\{\Pi_r\} \subset \mathcal{A}(P)\) such that:

\begin{enumerate}
\def\labelenumi{\arabic{enumi}.}
\tightlist
\item
  \textbf{Classicality:} \(\mathcal{R} := \text{vN}(\{\Pi_r\})\) is approximately abelian
\item
  \textbf{Stability:} \(\{\Pi_r\}\) are stable under effective dynamics
\item
  \textbf{Shareability:} \(\Pi_r\) lies in the center of overlap algebras
\end{enumerate}

\subsection{3.2 Markov Structure Theorem}\label{32-markov-structure-theorem}

\textbf{Theorem 3.1 (Quantum Markov chain structure on a fixed EC collar).} If \(I(A:C\mid B)=0\) for a tripartite state \(\rho_{ABC}\), then
\[
\mathcal H_B
\cong
\bigoplus_j
\left(
\mathcal H_{b_L^{(j)}}\otimes \mathcal H_{b_R^{(j)}}
\right)
\]
and the state decomposes as
\[
\rho_{ABC}
=
\bigoplus_j
q_j\,
\rho^{(j)}_{A b_L^{(j)}}\otimes \rho^{(j)}_{b_R^{(j)} C}.
\]
The label \(j\) is a classical variable living in the center of \(\mathcal A(B)\).

On the fixed finite-dimensional collar Hilbert space supplied by EC, let
\[
\mathfrak M_{A:B:C}
:=
\left\{
\sigma_{ABC}: I(A:C\mid B)_\sigma=0
\right\}
\]
and define the exact-Markov distance modulus
\[
\delta^{\mathrm M}_{A:B:C}(\varepsilon)
:=
\sup\left\{
\inf_{\sigma\in\mathfrak M_{A:B:C}}\|\rho-\sigma\|_1:
I(A:C\mid B)_\rho\le\varepsilon
\right\}.
\]
Because the state space is compact and conditional mutual information is continuous in finite dimension,
\[
\delta^{\mathrm M}_{A:B:C}(\varepsilon)\longrightarrow0
\qquad
(\varepsilon\downarrow0).
\]

This is the precise sense in which approximate collars converge to the exact HJPW normal form: either the state is exactly Markov, or one works on one fixed finite-dimensional collar, or a compatible pulled-back family, and sends \(\varepsilon\to0\). The supplement does \emph{not} claim a universal dimension-free one-shot trace-norm bound from small \(I(A:C\mid B)\) directly to an exact Markov state.

\subsection{3.3 Approximate Markov and Recovery}\label{33-approximate-markov-and-recovery}

\textbf{Theorem 3.2 (Fawzi--Renner recovery).} If
\[
I(A:C\mid B)_\rho\le\varepsilon,
\]
then there exists a CPTP map \(\mathcal R:\mathcal A(B)\to\mathcal A(BC)\) such that the recovered comparison state
\[
\rho^{\mathrm{rec}}_{ABC}
:=
(\mathrm{id}_A\otimes\mathcal R)(\rho_{AB})
\]
obeys
\[
\|\rho_{ABC}-\rho^{\mathrm{rec}}_{ABC}\|_1
\le
2\sqrt{1-e^{-\varepsilon}}
\le
2\sqrt{\varepsilon}.
\]
If \(\varepsilon\) is quoted in bits rather than nats, this bound is equivalently written as \(2\sqrt{\ln 2\,\varepsilon}\).

The recovered state is a constructive comparison state. It need not be exact Markov. To connect approximate recoverability to the exact splice or additivity formulas used later, one must combine Theorem 3.2 with the fixed-collar modulus of Section 3.2. Concretely, if \(\sigma_\varepsilon\in\mathfrak M_{A:B:C}\) is chosen so that
\[
\|\rho-\sigma_\varepsilon\|_1
\le
\delta^{\mathrm M}_{A:B:C}(\varepsilon),
\]
then:
\begin{enumerate}
\item every bounded observable supported on the left side of the collar differs from its exact-Markov splice value by at most \(\|X\|_\infty\,\delta^{\mathrm M}_{A:B:C}(\varepsilon)\);
\item if the relevant marginals are uniformly faithful with lower spectral bound \(\lambda_\ast>0\), then the corresponding modular-additivity defect differs from the exact Markov value by at most
\[
4\lambda_\ast^{-1}\,
\delta^{\mathrm M}_{A:B:C}(\varepsilon).
\]
\end{enumerate}

Thus the supplement carries two distinct errors:
\[
r_{\mathrm{FR}}(\varepsilon)=O(\varepsilon^{1/2})
\]
for constructive one-shot recovery, and
\[
\delta^{\mathrm M}_{A:B:C}(\varepsilon)\to0
\]
for convergence to the exact Markov normal form on a fixed collar. This is the controlled bridge from approximate recoverability to the exact splice and modular-additivity statements used by the later gravity and consensus arguments.

\subsection{3.4 Definiteness of Outcomes}\label{34-definiteness-of-outcomes}

\textbf{Theorem 3.3 (Collapse from EC Structure):} Under EC, the post-measurement physical state is exactly a convex mixture over sector labels \(\alpha\): \[\rho = \bigoplus_\alpha p_\alpha \, \rho^{(\alpha)}\]

For any observer whose accessible observables lie in either side algebra:

\begin{enumerate}
\def\labelenumi{\arabic{enumi}.}
\tightlist
\item
  There are \textbf{no interference observables} between different \(\alpha\)-blocks
\item
  Each observer-instance has a definite \(\alpha\) in the classical sense
\end{enumerate}

\emph{Proof:} Cross-block operators \(|i\rangle\langle j|\) for \(i \in \alpha, j \in \alpha'\) (\(\alpha \neq \alpha'\)) are not gauge-invariant, hence not in the physical algebra. \(\square\)

\subsection{3.5 Born Rule as Unique Probability Measure}\label{35-born-rule-as-unique-probability-measure}

\textbf{Theorem 3.4 (Gleason):} For \(\dim \mathcal{H} \geq 3\), any measure \(\mu\) on projectors satisfying:

\begin{enumerate}
\def\labelenumi{\arabic{enumi}.}
\tightlist
\item
  Noncontextuality: \(\mu(P)\) depends only on \(P\)
\item
  Additivity: \(\mu(P \vee Q) = \mu(P) + \mu(Q)\) for \(P \perp Q\)
\item
  Normalization: \(\mu(\mathbf{1}) = 1\)
\end{enumerate}

must have the form: \[\mu(P) = \text{Tr}(\rho P)\] for some density operator \(\rho\).

\textbf{Corollary 3.5:} In OPH with the type-I regulator premise (R0), Born\textquotesingle s rule is the unique overlap-consistent probability assignment.

\subsection{3.6 Collapse as Conditioning}\label{36-collapse-as-conditioning}

\textbf{Proposition 3.6 (L\"{u}ders Update):} For record projectors \(\{\Pi_r\}\) and state \(\omega\): \[p_r = \omega(\Pi_r), \qquad \omega_r(X) = \frac{\omega(\Pi_r X \Pi_r)}{\omega(\Pi_r)}\]

Then for all \(X\) commuting with \(\mathcal{R}\): \[\omega(X) = \sum_r p_r \, \omega_r(X)\]

This is classical conditional probability on the record algebra.

\begin{center}\rule{0.5\linewidth}{0.5pt}\end{center}

\section{4. The Cosmological Principle}\label{4-the-cosmological-principle}

\subsection{4.1 MaxEnt Symmetry Inheritance}\label{41-maxent-symmetry-inheritance}

\textbf{Lemma 4.1:} Let \(G\) act on the screen by automorphisms \(\alpha_g\). If:

\begin{enumerate}
\def\labelenumi{\arabic{enumi}.}
\tightlist
\item
  The constraint set is \(G\)-invariant
\item
  Von Neumann entropy is \(G\)-invariant: \(S(\alpha_g(\rho)) = S(\rho)\)
\end{enumerate}

Then the MaxEnt optimizer \(\omega\) can be taken \(G\)-invariant. If unique, it is automatically \(G\)-invariant.

\emph{Proof:} If \(\omega\) is a maximizer, so is \(\omega \circ \alpha_{g^{-1}}\) by constraint invariance and entropy invariance. By uniqueness, \(\omega = \omega \circ \alpha_{g^{-1}}\). \(\square\)

\subsection{4.2 Isotropy from SO(3)-Invariant Constraints}\label{42-isotropy-from-so3-invariant-constraints}

\textbf{Assumption C:} Constraints are SO(3)-invariant on \(S^2\).

\textbf{Theorem 4.2:} Under Assumption C and MaxEnt uniqueness, the reference state \(\omega\) satisfies: \[\omega \circ \alpha_g = \omega \quad \forall g \in \text{SO}(3)\]

The stress tensor has perfect-fluid form: \[\langle T_{ab} \rangle = \rho \, u_a u_b + p(g_{ab} + u_a u_b)\]

\emph{Proof of perfect-fluid form:} Under SO(3), rank-2 tensors decompose into scalar + vector + traceless tensor irreps. Only scalars are invariant. Hence \(\langle T_{0i} \rangle = 0\) and \(\langle T_{ij} \rangle \propto \delta_{ij}\). \(\square\)

\subsection{4.3 Homogeneity from Isotropy Everywhere}\label{43-homogeneity-from-isotropy-everywhere}

\textbf{Theorem 4.3 (Schur-type):} Let \((\Sigma, h_{ij})\) be a 3-dimensional Riemannian manifold. If at every point \(p \in \Sigma\), the curvature tensor is SO(3)-invariant (pointwise isotropy), then: \[R_{ijkl}(p) = K(p)(h_{ik}h_{jl} - h_{il}h_{jk})\]

and by the second Bianchi identity, \(\nabla_m K = 0\), so \(K = \text{const}\).

\textbf{Corollary 4.4:} If OPH supplies isotropy for all observers, spatial geometry is a constant-curvature space form: \(S^3\), \(\mathbb{R}^3\), or \(H^3\).

\subsection{4.4 FLRW Emergence}\label{44-flrw-emergence}

\textbf{Theorem 4.5:} With:

\begin{enumerate}
\def\labelenumi{\arabic{enumi}.}
\tightlist
\item
  Semiclassical Einstein equation from recoverable generalized entropy / entanglement equilibrium
\item
  Perfect-fluid stress tensor from SO(3)-invariant MaxEnt
\item
  Positive \(\Lambda\) from finite screen capacity
\end{enumerate}

The emergent geometry is FLRW: \[ds^2 = -dt^2 + a(t)^2 \left(\frac{dr^2}{1-kr^2} + r^2 d\Omega^2\right)\]

with Friedmann equations: \[H^2 + \frac{k}{a^2} = \frac{8\pi G}{3}\rho + \frac{\Lambda}{3}\] \[\dot{H} - \frac{k}{a^2} = -4\pi G(\rho + p)\]

\begin{center}\rule{0.5\linewidth}{0.5pt}\end{center}

\section{5. Supersymmetry and Gauge Coupling Unification}\label{5-supersymmetry-and-gauge-coupling-unification}

\subsection{5.1 One-Loop Unification Algebra}\label{51-one-loop-unification-algebra}

At one loop, if couplings unify at \((M_U, \alpha_U)\): \[\alpha_i^{-1}(M_Z) = \alpha_U^{-1} + \frac{b_i}{2\pi}\ln\frac{M_U}{M_Z}\]

Define \(A_i := \alpha_i^{-1}(M_Z)\), \(L := \ln(M_U/M_Z)\). Eliminating \(\alpha_U\): \[L = \frac{2\pi}{b_1 - b_2}(A_1 - A_2)\]

The consistency relation for \(\alpha_s\) (a consistency check when all three couplings calibrate $P$; a genuine prediction in the two-input mode): \[A_3^{\text{pred}} = \frac{b_3 - b_2}{b_1 - b_2}A_1 + \frac{b_1 - b_3}{b_1 - b_2}A_2\]

\subsection{5.2 SM vs MSSM Coefficients}\label{52-sm-vs-mssm-coefficients}

{\def\LTcaptype{none} % do not increment counter
\begin{longtable}[]{@{}>{\RaggedRight\arraybackslash}p{0.320\textwidth}>{\RaggedRight\arraybackslash}p{0.320\textwidth}>{\RaggedRight\arraybackslash}p{0.320\textwidth}@{}}
\toprule\noalign{}
Model & \((b_1, b_2, b_3)\) & \(\alpha_s(M_Z)^{\text{pred}}\) \\
\midrule\noalign{}
\endhead
\bottomrule\noalign{}
\endlastfoot
SM & \((41/10, -19/6, -7)\) & \(\approx 0.071\) \\
MSSM & \((33/5, 1, -3)\) & \(\approx 0.116\) \\
Observed & , & \(0.1179 \pm 0.0010\) \\
\end{longtable}
}

The SM prediction is catastrophically wrong; MSSM-like coefficients work.

\subsection{5.3 Heat-Kernel Edge Sector Weights}\label{53-heat-kernel-edge-sector-weights}

\textbf{Theorem 5.1:} Under MaxEnt with bi-invariant constraints on a compact Lie group \(G\), the sector probabilities take heat-kernel form: \[p_R(t) \propto d_R \, e^{-t \, C_2(R)}\]

where \(d_R = \dim R\) and \(C_2(R)\) is the quadratic Casimir.

\textbf{Lemma 5.2 (Bi-invariant Operators):} If \(G=\prod_i G_i\) is a compact semisimple Lie group written as a product of compact simple factors, then any bi-invariant, second-order differential operator on \(G\) has the form: \[D = c_0 \mathbf{1} - \sum_i c_i \Delta_{G_i}\]

where \(\Delta_G\) is the Laplace-Beltrami operator.

\emph{Proof:} Bi-invariance \(\Leftrightarrow\) \(D \in Z(U(\mathfrak{g}))\) (center of universal enveloping algebra). For degree \(\leq 2\):

\begin{itemize}
\tightlist
\item
  Linear terms vanish (no invariant vectors in adjoint rep)
\item
  Quadratic terms are proportional to the Killing form on each simple factor, giving one independent coefficient per factor
\end{itemize}

The Casimir element acts as \(-\Delta_G\) in the regular representation. \(\square\)

\subsection{5.4 Peter-Weyl and Effective Multiplicity}\label{54-peter-weyl-and-effective-multiplicity}

By Peter-Weyl decomposition, the effective refinement-limit edge representation space is
\[L^2(G) \cong \bigoplus_R V_R \otimes V_R^*\]

\textbf{Key insight:} Entanglement traces over one factor give multiplicity \(d_R\) in \(p_R\), but loops see both factors, giving effective multiplicity: \[N_{\text{eff}}(R) = d_R \cdot p_R\]

\subsection{5.5 Beta Function Shifts}\label{55-beta-function-shifts}

The one-loop beta shift from edge sectors: \[\Delta b_a = \sum_R p_R \cdot d_R \cdot T_a(R)\]

where \(T_a(R)\) is the Dynkin index for gauge factor \(a\).

At the unification diffusion parameter \(t_U \approx 1.64\): \[\Delta b \approx (2.49, 4.38, 3.97) \quad \text{vs MSSM} \quad (2.50, 4.17, 4.00)\]

With fermionic grading (restricting to half-integer SU(2) sectors): \[\Delta b \approx (2.50, 4.17, 3.97)\]

matching MSSM within 1\%.

\begin{center}\rule{0.5\linewidth}{0.5pt}\end{center}

\section{6. The Horizon Problem}\label{6-the-horizon-problem}

\subsection{6.1 Anisotropy Bound from Markov Control}\label{61-anisotropy-bound-from-markov-control}

\textbf{Theorem 6.1:} For a tripartition \(A\)-\(B\)-\(D\) with collar width \(\delta\), the Markov error satisfies: \[I(A:D|B) \lesssim c \cdot |\partial C|_{\text{UV}} \cdot e^{-\delta/\xi}\]

where \(\xi\) is the correlation length.

\textbf{Corollary 6.2:} For bounded observable \(O\) with \(\|O\| \sim 1\): \[|\Delta\langle O \rangle| \leq 2\sqrt{\ln 2 \cdot \varepsilon}\]

\subsection{6.2 CMB Anisotropy Requirement}\label{62-cmb-anisotropy-requirement}

For \(\delta T/T \lesssim 10^{-5}\): \[2\sqrt{\ln 2 \cdot \varepsilon} \lesssim 10^{-5}\] \[\varepsilon \lesssim 3.61 \times 10^{-11} \text{ bits}\]

Required collar width: \[\frac{\delta}{\xi} \gtrsim \ln\left(\frac{c \cdot |\partial C|_{\text{UV}}}{\varepsilon_{\text{max}}}\right) \approx 24\]

With \(\xi \approx \sqrt{a_{\text{cell}}} \cdot \ell_P \approx 1.28 \ell_P\): \[\delta_{\text{CMB}} \approx 31 \ell_P\]

\subsection{6.3 Homogeneity as MaxEnt Default}\label{63-homogeneity-as-maxent-default}

\textbf{Theorem 6.3:} If MaxEnt constraints are uniform across UV cells with no marked points, the MaxEnt state is invariant under triangulation automorphisms. In the continuum limit, this gives SO(3) invariance.

\emph{Proof:} Standard MaxEnt uniqueness + symmetry inheritance (Lemma 4.1). \(\square\)

\begin{center}\rule{0.5\linewidth}{0.5pt}\end{center}

\section{7. Black Hole Information Paradox}\label{7-black-hole-information-paradox}

\subsection{7.1 The Information Trilemma}\label{71-the-information-trilemma}

The AMPS-style paradox assumes:

\begin{enumerate}
\def\labelenumi{\arabic{enumi}.}
\tightlist
\item
  Outgoing mode \(B\) entangled with interior partner \(A\) (smooth horizon)
\item
  Late radiation \(B\) entangled with early radiation \(R\) (unitarity)
\item
  Monogamy: \(B\) cannot be maximally entangled with both
\end{enumerate}

\textbf{The false assumption:} \(\mathcal{H} \stackrel{?}{=} \mathcal{H}_{\text{inside}} \otimes \mathcal{H}_{\text{outside}} \otimes \mathcal{H}_{\text{radiation}}\)

\subsection{7.2 Edge-Center Resolution}\label{72-edge-center-resolution}

\textbf{Theorem 7.1:} Under EC, the collar decomposition gives: \[\rho_{A_\delta B_\delta D_\delta} = \bigoplus_\alpha p_\alpha \left(\rho_{A_\delta b_L^\alpha} \otimes \rho_{b_R^\alpha D_\delta}\right)\]

The "glue" between inside and outside is the edge sector label \(\alpha\) in the center. Conditional on \(\alpha\), inside and outside factorize.

\subsection{7.3 Recoverability of Interior}\label{73-recoverability-of-interior}

By the recoverability clause of the canonical entropy axiom (small CMI) + Theorem 3.2: \[\left\|\rho_{ABC} - (\text{id}_A \otimes \mathcal{R})(\rho_{AB})\right\|_1 \leq 2\sqrt{\ln 2 \cdot \varepsilon}\]

\textbf{Interpretation:} The "interior partner" is recoverable from outside + radiation. It\textquotesingle s encoded, not independent.

\subsection{7.4 Hawking Temperature from Modular Theory}\label{74-hawking-temperature-from-modular-theory}

\textbf{Theorem 7.2:} The outside algebra is in a KMS state at inverse temperature \(\beta = 2\pi/\kappa\) with respect to the horizon Killing generator.

\emph{Derivation:} Near the horizon: \[ds^2 \approx -\kappa^2 \rho^2 dt^2 + d\rho^2 + r_H^2 d\Omega^2\]

Euclidean regularity at \(\rho = 0\) requires \(\kappa t_E \sim \kappa t_E + 2\pi\), hence: \[T_H = \frac{\kappa}{2\pi} = \frac{\hbar c^3}{8\pi G k_B M}\]

\subsection{7.5 Discrete Area Spectrum}\label{75-discrete-area-spectrum}

From the central area operator: \[L_C = \sum_\alpha (\log d_\alpha) P_\alpha\]

Area eigenvalues: \(A_\alpha = 4G \log d_\alpha = 4\ell_P^2 \ln d_\alpha\)

For Schwarzschild, a sector transition \(d \to d'\) gives: \[\Delta E = k_B T_H \ln(d'/d)\]

\begin{center}\rule{0.5\linewidth}{0.5pt}\end{center}

\section{8. The Cosmological Constant Problem}\label{8-the-cosmological-constant-problem}

\subsection{8.1 Vacuum Energy Blindness}\label{81-vacuum-energy-blindness}

\textbf{Lemma 8.1:} For any null vector \(k\): \[T^{\text{vac}}_{kk} = T^{\text{vac}}_{ab} k^a k^b = -\rho_{\text{vac}} g_{ab} k^a k^b = -\rho_{\text{vac}}(k \cdot k) = 0\]

Vacuum energy contributions lie in the kernel of the null map \(T_{ab} \mapsto T_{kk}\).

\subsection{8.2 Structural Separation}\label{82-structural-separation}

\textbf{Proposition 8.2:} In OPH:

\begin{itemize}
\tightlist
\item
  Local modular/null data fixes \(T_{ab}\) up to \(\phi g_{ab}\)
\item
  \(\Lambda\) is fixed by global screen capacity, not local physics
\end{itemize}

The "huge zero-point energy" simply isn\textquotesingle t the thing determining curvature.

\subsection{8.3 $\Lambda$ from Capacity}\label{83-ux3bb-from-capacity}

\[\Lambda = \frac{3\pi}{G \cdot \log(\dim \mathcal{H}_{\text{tot}})}\]

Equivalently, using de Sitter entropy: \[S_{dS} = \frac{A_{dS}}{4G} = \frac{3\pi}{G\Lambda} = \log(\dim \mathcal{H}_{\text{tot}})\]

For observed \(\Lambda \approx 1.09 \times 10^{-52}\) m\(^{-2}\): \[\log(\dim \mathcal{H}_{\text{tot}}) \sim 10^{122}\]

\subsection{8.4 No-Go for Local $\Lambda$ Prediction}\label{84-no-go-for-local-ux3bb-prediction}

\textbf{Theorem 8.3:} Within OPH\textquotesingle s null-modular reconstruction:

\begin{enumerate}
\def\labelenumi{\arabic{enumi}.}
\tightlist
\item
  All \(T_{kk}\) data unchanged under \(T_{ab} \to T_{ab} + \phi g_{ab}\)
\item
  Therefore \(\Lambda\) cannot be fixed by local overlap consistency
\item
  \(\Lambda\) requires a global input (capacity)
\end{enumerate}

\begin{center}\rule{0.5\linewidth}{0.5pt}\end{center}

\section{9. Dark Matter as Modular Anomaly}\label{9-dark-matter-as-modular-anomaly}

\subsection{9.1 The Modular Additivity Defect}\label{91-the-modular-additivity-defect}

Define the collar modular additivity defect: \[\Delta K_\delta := K_{ABD} - K_{AB} - K_{BD} + K_B\]

\textbf{Theorem 9.1:} \(\langle \Delta K_\delta \rangle_\omega = -I(A:D|B)_\omega\)

This is the conditional mutual information, measuring Markov imperfection.

\subsection{9.2 Anomalous Stress-Energy}\label{92-anomalous-stress-energy}

The modified Einstein equation: \[G_{ab} + \Lambda g_{ab} = 8\pi G \left(\langle T_{ab} \rangle + \langle T_{ab}^{\text{anom}} \rangle\right)\]

where in the diamond rest frame: \[\langle T_{00}^{\text{anom}} \rangle := \frac{15}{8\pi^2} \frac{\delta\langle K_C^{(\text{anom})} \rangle}{\ell^4}\]

The coefficient \(15/(8\pi^2)\) comes from inverting the \(d=4\) geometric integral coefficient \(\Omega_2/(4^2-1) = 4\pi/15\).

\subsection{9.3 Properties of the Anomaly Sector}\label{93-properties-of-the-anomaly-sector}

\begin{enumerate}
\def\labelenumi{\arabic{enumi}.}
\tightlist
\item
  \textbf{Gravitates:} Appears on RHS of Einstein equation
\item
  \textbf{Dark by construction:} Encoded in central/recoverability data, not SM fields
\item
  \textbf{Covariantly conserved:} \(\nabla^a T_{ab}^{\text{anom}} = 0\) (from Bianchi identity)
\item
  \textbf{Effectively classical:} Central structure implies classical labels
\end{enumerate}

\subsection{9.4 The MOND Acceleration Scale}\label{94-the-mond-acceleration-scale}

\textbf{Corollary 9.2 (program-level):} The unique IR acceleration scale available to the modular-anomaly branch from \(\Lambda\) is \[a_0^{(\text{OPH})} := \frac{15}{8\pi^2} c^2 \sqrt{\frac{\Lambda}{3}} = \frac{15}{8\pi^2} \frac{c^2}{r_{dS}}\]

\textbf{Numerical evaluation:} \[a_0^{(\text{OPH})} \approx 1.03 \times 10^{-10} \text{ m/s}^2\]

Compare to observed: \(a_0^{(\text{obs})} \approx 1.2 \times 10^{-10}\) m/s\(^2\) (within 15\%).

\subsection{9.5 Polarization Response (MOND-like Scaling)}\label{95-polarization-response-mond-like-scaling}

\textbf{Uniqueness argument:} In the deep-IR regime, the only scales are \(a_0\) and \(g_b\) (Newtonian baryonic acceleration). Dimensional analysis + scale invariance forces: \[g_{\text{anom}} \propto \sqrt{a_0 \cdot g_b}\]

The cubic polarization functional: \[S_{\text{pol}}[\varphi] = \int d^3x \left[-\frac{1}{12\pi G a_0}|\nabla\varphi|^3 - \rho_b \varphi\right]\]

yields the MOND/RAR form: \[g_{\text{obs}} \approx g_b + \sqrt{a_0 g_b}\]

\begin{center}\rule{0.5\linewidth}{0.5pt}\end{center}

\section{10. Proton Stability and Spin}\label{10-proton-stability-and-spin}

This section separates three distinct claim types that were too compressed in older drafts:
\begin{enumerate}
\item \textbf{Structural output already closed on the realized D8--D9 branch:} the realized low-energy gauge group is a product up to finite quotient, so the adjoint contains no mixed \(X,Y\)-type gauge bosons and therefore no \emph{gauge-mediated} proton-decay channel of simple-GUT type.
\item \textbf{QCD/EFT-dependent hadron claims:} any finite proton-lifetime estimate beyond that gauge-channel exclusion, and any statement about the proton spin decomposition, require explicit operator matching, renormalization conventions, and nonperturbative hadronic matrix elements.
\item \textbf{Still-open completion burden:} the present supplement does not yet contain the nonperturbative light-quark/hadron completion needed to turn those QCD-dependent statements into derived OPH outputs.
\end{enumerate}

\subsection{10.1 Sector Factorization and Product Group}\label{101-sector-factorization-and-product-group}

\textbf{Theorem 10.1 (Factorization Equivalence):} \[\text{Factorizing edge weights} \Leftrightarrow \text{Additive boundary Laplacian} \Leftrightarrow \text{Product gauge group}\]

If \(H_\partial = H_\partial^{(1)} + H_\partial^{(2)} + H_\partial^{(3)}\) with \([H_\partial^{(i)}, H_\partial^{(j)}] = 0\): \[p(R_1, R_2, R_3) \propto \prod_{i=1}^3 d_{R_i} e^{-t_i C_2(R_i)}\]

Applied to the refinement-stable edge-sector colimit, Tannaka-Krein reconstruction first yields some compact group \(G\). Under the extended MAR admissibility package, and once the admissible class includes one connected abelian charge factor, the connected Lie realization selected by MAR is \(SU(3) \times SU(2) \times U(1)\) (up to finite quotient).

Under the extended theory \(T_{\mathrm{ext}}\), the explicit MAR selector picks this product structure on the admissible class: the minimal coupled carrier \(\mathbb{C}^3 \otimes \mathbb{C}^2\) enforces commuting color and weak actions, which implies the additive boundary Laplacian. See Part II of this manuscript for the complete proof.

\subsection{10.2 No Leptoquarks, No Gauge Proton Decay}\label{102-no-leptoquarks-no-gauge-proton-decay}

\textbf{Corollary 10.2 (structural D8--D9 consequence only):} With product gauge group, the adjoint representation of the full connected gauge group is \[(8,1,0) \oplus (1,3,0) \oplus (1,1,0)\]

equivalently, the adjoint of the derived gauge group is
\[
(8,1,0) \oplus (1,3,0).
\]

There are therefore no gauge generators in mixed representations \((3,2,\pm 5/6)\), i.e.\ no simple-GUT \(X,Y\) bosons on the realized OPH branch. The structural claim is exactly
\[
\tau_p^{(\text{gauge})} = \infty,
\]
meaning that the gauge-boson proton-decay channel is absent.

\textbf{What is not proved here.} This does \emph{not} establish full proton stability, and it does not justify the stronger shorthand \(\tau_p=\infty\). Scalar-mediated or higher-dimensional baryon-violating operators, RG transport of those operators from a UV completion, and the hadronic matrix elements entering exclusive proton-decay rates are all outside the present structural argument. Those steps are EFT/QCD-dependent and remain open in the current supplement. The electroweak baryon-violating counting discussed in Section~\ref{11-baryogenesis} is a separate high-temperature/topological issue rather than a zero-temperature proton-stability theorem.

\subsection{10.3 Proton Spin Fraction}\label{103-proton-spin-fraction}

\textbf{Structural content already fixed.} On the realized D8--D9 branch the color factor is \(\mathrm{SU}(3)\), so the relevant Casimirs are
\[
C_F=\frac43,\qquad C_A=3.
\]
That is the full structural OPH input currently available for proton-spin bookkeeping.

\textbf{QCD-dependent continuation ansatz.} If one adds the phenomenological ansatz that quark/gluon spin sharing equilibrates according to the color Casimirs, then
\[
\Delta\Sigma \approx \frac{C_F}{C_F + C_A}.
\]

For \(\mathrm{SU}(3)\) this gives
\[
\Delta\Sigma = \frac{4/3}{4/3 + 3} = \frac{4}{13} \approx 0.308.
\]

Compared to the representative lattice benchmark \(\Delta\Sigma \approx 0.286\), this lands within about \(8\%\). But the displayed value is only a deferred continuation benchmark, not a theorem. The observable \(\Delta\Sigma\) is factorization-scheme and scale dependent and, in QCD, depends on the polarized axial-current matrix element together with gluon-spin and orbital-angular-momentum contributions.

\textbf{Still open.} A first-principles OPH derivation of proton spin would require the nonperturbative light-quark/hadron completion of Phase I plus an explicit map from OPH data to renormalized proton matrix elements. Until that exists, the \(4/13\) estimate should be read only as a QCD-dependent continuation ansatz.

\begin{center}\rule{0.5\linewidth}{0.5pt}\end{center}

\section{11. Baryogenesis}\label{11-baryogenesis}

\subsection{11.1 The Z$_6$ Defect Suppression}\label{111-the-zux2086-defect-suppression}

The SM global gauge structure is \((G_{\text{SM}})/\mathbb{Z}_6\). The entropy deficit from the \(\mathbb{Z}_6\) quotient: \[\Delta S = \ln 6\]

Defect suppression factor: \[\varepsilon = e^{-\Delta S} = \frac{1}{6}\]

\subsection{11.2 The Baryon-Violating Channel}\label{112-the-baryon-violating-channel}

The electroweak topological transition (\(\Delta N_{CS} = \pm 1\)) produces chiral zero modes. By index theorem:

\begin{itemize}
\tightlist
\item
  One zero mode per left-handed SU(2) doublet per unit topological charge
\item
  Per generation: \(N_c + 1 = 4\) doublets
\item
  Total: \(n_{\text{doublets}} = (N_c + 1)N_g = 12\)
\end{itemize}

The \textquotesingle t Hooft vertex is a 12-fermion interaction.

\subsection{11.3 Defect-Order Identification}\label{113-defect-order-identification}

\textbf{Proposition 11.1:} The baryon-violating channel corresponds to defect order \(n = 12\): \[\varepsilon^{12} = 6^{-12} \approx 4.59 \times 10^{-10}\]

Compare to observed \(\eta_B \approx 6.1 \times 10^{-10}\) (same order of magnitude).

\subsection{11.4 CP Bias from Overlap Holonomy}\label{114-cp-bias-from-overlap-holonomy}

Under CP: \(U_{ij} \mapsto U_{ij}^*\), \(\Omega_{ijk} \mapsto \Omega_{ijk}^\dagger\)

For central cocycle \(z_{ijk} \in U(1)\): CP sends \(z \to z^{-1}\).

\textbf{Real cocycle condition:} \([z] = [z]^{-1}\) (cohomologous to inverse)

Generic phases violate this \(\Rightarrow\) CP violation is generic in OPH.

\begin{center}\rule{0.5\linewidth}{0.5pt}\end{center}

\section{12. Generation Structure and Yukawa Hierarchy}\label{12-generation-structure-and-yukawa-hierarchy}

\subsection{12.1 Why Three Generations}\label{121-why-three-generations}

\textbf{Theorem 12.1 (N\_g = 3 Selection):}

\begin{enumerate}
\def\labelenumi{\arabic{enumi}.}
\tightlist
\item
  \textbf{CP violation requires N\_g ≥ 3:} Physical CKM phase exists only for \(N_g \geq 3\)
\item
  \textbf{UV stability bounds N\_g ≤ 5:} SU(2) asymptotic freedom
\item
  \textbf{Minimality selects N\_g = 3:} Extra generations introduce unconstrained relevant deformations
\end{enumerate}

\subsection{12.2 Yukawa as Defect-Mediated Overlap}\label{122-yukawa-as-defect-mediated-overlap}

\textbf{Theorem 12.2:} If a Yukawa coupling requires \(n\) units of \(\mathbb{Z}_6\) defect insertion: \[y \propto \langle D^n \rangle = 6^{-n}\]

The suppression factor is fixed by topology, not tuned.

\subsection{12.3 The Yukawa Spectrum}\label{123-the-yukawa-spectrum}

{\def\LTcaptype{none} % do not increment counter
\begin{longtable}[]{@{}>{\RaggedRight\arraybackslash}p{0.240\textwidth}>{\RaggedRight\arraybackslash}p{0.240\textwidth}>{\RaggedRight\arraybackslash}p{0.240\textwidth}>{\RaggedRight\arraybackslash}p{0.240\textwidth}@{}}
\toprule\noalign{}
Fermion & Yukawa & \(n = -\ln y / \ln 6\) & Nearest integer \\
\midrule\noalign{}
\endhead
\bottomrule\noalign{}
\endlastfoot
\(t\) & 0.992 & 0.00 & 0 \\
\(b\) & 0.024 & 2.08 & 2 \\
\(c\) & 0.0073 & 2.75 & 3 \\
\(s\) & 5.3×10⁻⁴ & 4.21 & 4 \\
\(d\) & 2.7×10⁻⁵ & 5.87 & 6 \\
\(u\) & 1.3×10⁻⁵ & 6.29 & 6 \\
\(\tau\) & 0.010 & 2.56 & 3 \\
\(\mu\) & 6.1×10⁻⁴ & 4.13 & 4 \\
\(e\) & 2.9×10⁻⁶ & 7.11 & 7 \\
\end{longtable}
}

The base-6 logarithm lands near integers across the spectrum.

\subsection{12.4 Deterministic Completion (RK Hamiltonian)}\label{124-deterministic-completion-rk-hamiltonian}

Define configuration weight: \[w(\mathcal{C}) = \exp\left(-t \sum_{\text{plaquettes}} C_2(R_p)\right) \times 6^{-n_D(\mathcal{C})}\]

The Rokhsar-Kivelson ground state: \[|\Psi\rangle = \frac{1}{\sqrt{Z}} \sum_{\mathcal{C}} \sqrt{w(\mathcal{C})} |\mathcal{C}\rangle\]

is the unique ground state of a frustration-free parent Hamiltonian, replacing MaxEnt ensemble with a deterministic pure state.

\begin{center}\rule{0.5\linewidth}{0.5pt}\end{center}

\section{13. The Koide Formula}\label{13-the-koide-formula}

\subsection{13.1 Numerical Verification}\label{131-numerical-verification}

Using PDG masses:

\begin{itemize}
\tightlist
\item
  \(m_e = 0.51099895\) MeV
\item
  \(m_\mu = 105.6583755\) MeV
\item
  \(m_\tau = 1776.93\) MeV
\end{itemize}

\[Q = \frac{m_e + m_\mu + m_\tau}{(\sqrt{m_e} + \sqrt{m_\mu} + \sqrt{m_\tau})^2} = 0.6666644645\]

Compare to \(2/3 = 0.6666666667\). Difference: \(2.2 \times 10^{-6}\) (within 1σ).

\subsection{13.2 Z$_3$ Holonomy Parameterization}\label{132-zux2083-holonomy-parameterization}

\textbf{Theorem 13.1:} The minimal Hermitian circulant on generation space: \[\Phi = a \cdot I + b \cdot P + b^* \cdot P^2\]

with \(P\) the cyclic shift, has eigenvalues: \[\lambda_k = a + 2|b|\cos\left(\delta + \frac{2\pi k}{3}\right)\]

where \(\delta = \arg(b)\).

\subsection{13.3 Why Q = 2/3}\label{133-why-q--23}

\textbf{Theorem 13.2:} With \(r_k = \lambda_k\) (root masses): \[Q = \frac{\sum_k r_k^2}{(\sum_k r_k)^2} = \frac{1 + 2(|b|/a)^2}{3}\]

Setting \(Q = 2/3\): \[|b|/a = \frac{1}{\sqrt{2}}\]

This is the "balanced" configuration where singlet and charged modes have equal norm.

\subsection{13.4 The Koide Phase from OPH}\label{134-the-koide-phase-from-oph}

\textbf{Proposition 13.3:} The holonomy phase: \[\delta_{\text{OPH}} = \frac{\beta_{\text{EW}} \cdot Y_Q}{N_g} = \frac{(N_c + 1)}{2N_c \cdot N_g}\]

With \(N_c = 3\), \(N_g = 3\): \[\delta_{\text{OPH}} = \frac{4}{18} = \frac{2}{9} = 0.2222...\]

Experimental extraction: \(\delta_{\text{exp}} = 0.2222248 \pm 0.0000063\)

Agreement within 0.4σ.

\begin{center}\rule{0.5\linewidth}{0.5pt}\end{center}

\section{14. Cosmological Parameters}\label{14-cosmological-parameters}

\subsection{14.1 The de Sitter Scale}\label{141-the-de-sitter-scale}

From \(\Lambda \approx 1.09 \times 10^{-52}\) m\(^{-2}\): \[r_{dS} = \sqrt{\frac{3}{\Lambda}} \approx 1.66 \times 10^{26} \text{ m}\]

The de Sitter time: \[t_\Lambda = \frac{r_{dS}}{c} \approx 5.53 \times 10^{17} \text{ s} \approx 17.5 \text{ Gyr}\]

\subsection{14.2 Emergent Age of the Universe}\label{142-emergent-age-of-the-universe}

For flat ΛCDM: \[t_0 = \frac{2}{3H_0\sqrt{\Omega_\Lambda}} \sinh^{-1}\left(\sqrt{\frac{\Omega_\Lambda}{\Omega_m}}\right)\]

With \(H_0 \approx 67.4\) km/s/Mpc, \(\Omega_\Lambda \approx 0.685\): \[t_0 \approx 13.8 \text{ Gyr}\]

\subsection{14.3 Screen Capacity}\label{143-screen-capacity}

\[\dim \mathcal{H}_{\text{tot}} = \exp(S_{dS}) = \exp\left(\frac{3\pi}{G\Lambda}\right)\]

\[\log(\dim \mathcal{H}_{\text{tot}}) \approx 2.85 \times 10^{122}\]

\subsection{14.4 Pixel Parameters}\label{144-pixel-parameters}

\begin{itemize}
\tightlist
\item
  Cell area: \(a_{\text{cell}} \approx 1.63 \ell_P^2\)
\item
  UV length: \(\ell_{\text{UV}} = \sqrt{a_{\text{cell}}} \approx 1.28 \ell_P\)
\item
  Entropy per cell: \(\bar{\ell} = a_{\text{cell}}/(4\ell_P^2) \approx 0.408\)
\end{itemize}

\begin{center}\rule{0.5\linewidth}{0.5pt}\end{center}

\section{Appendix A: Summary of Derived Results}\label{appendix-a-summary-of-derived-results}

\subsection{A.1 Theorem-Level Results and Dependency Status}\label{a1-theorem-level-results}

This appendix records theorem-level status only. The recovered core of the current OPH program is D1--D5 together with D7--D9. D6 is included below only as the adjacent Phase-II boundary marker; D10--D12 are excluded from the theorem-level list except where their non-core status must be made explicit. The table matches the compact note's authoritative ledger and keeps imported mathematics, extra premises, and claim tier visible.

\begingroup
\scriptsize
{\def\LTcaptype{none}
\setlength{\tabcolsep}{2pt}
\begin{longtable}[]{@{}>{\RaggedRight\arraybackslash}p{0.07\textwidth}>{\RaggedRight\arraybackslash}p{0.21\textwidth}>{\RaggedRight\arraybackslash}p{0.15\textwidth}>{\RaggedRight\arraybackslash}p{0.22\textwidth}>{\RaggedRight\arraybackslash}p{0.17\textwidth}@{}}
\toprule\noalign{}
Node & Output & Standard mathematics used & Extra premises / inputs beyond the five axioms & Status \\
\midrule\noalign{}
\endhead
\bottomrule\noalign{}
\endlastfoot
\textbf{D1} & schedule-independent normal form for overlap repair & Newman's lemma & finite patch net, termination, and local confluence & Phase I structural theorem \\
\textbf{D2} & EC decomposition, Markov collar, and generalized-entropy split & HJPW; Petz; Fawzi--Renner & R0/R1 or explicit recoverability control; coarse-grained \(A/(4G)\) dictionary & Phase I structural lemma/proposition layer \\
\textbf{D3} & geometric modular flow and Lorentz branch & Bisognano--Wichmann modular geometry; conformal-group classification & controlled refinement limit together with the OPH geometric-branch hypothesis of Theorem 2.2 & Phase I conditional scaling-limit theorem \\
\textbf{D4} & null modular bridge to \(T_{kk}\) & Borchers--Wiesbrock; distributional differentiation on half-lines & the theorem-local inherited-strip decomposition and exact-or-controlled Markov hypotheses of Section~\ref{23-null-modular-bridge-eft-conditions-n1-n3}, plus the downstream density-upgrade and relativistic null-stress-identification premises; the endpoint/propagation controls are internal to the local finite-constraint MaxEnt branch & Phase I conditional bridge theorem \\
\textbf{D5} & Jacobson-type rest-frame relation plus the conditional tensor upgrade & Jacobson small-ball mathematics; null-to-tensor reconstruction modulo a metric term & fixed-cap stationarity, the locally Lorentzian \(d=4\) scaling regime, and the all-directions/all-reference-states premise used in the tensor upgrade & Phase I conditional scaling-limit theorem/corollary \\
\textbf{D6} & capacity/\(\Lambda\) corollary and the separate capacity-based neutrino estimate & de Sitter entropy relation and dimensional analysis & external input \(N_{\mathrm{scr}}\) and the screen-capacity identification & Phase II input-dependent corollary / estimate \\
\textbf{D7} & compact gauge reconstruction in the bosonic branch & Doplicher--Roberts / Tannaka reconstruction & bosonic symmetric transportable colimit and fiber-functor premises; the refinement-stable qualifier is internal to the local finite-constraint MaxEnt/refinement branch & Phase I conditional structural theorem \\
\textbf{D8} & product structure up to finite quotient on the realized branch & compact Lie representation classification; Schur's lemma & MAR, connected admissible class, one connected abelian factor, and \([z]=0\) & Phase I realized-branch theorem \\
\textbf{D9} & exact hypercharges, the counting chain \(N_g=3\) then \(N_c=3\), and the \(\mathbb Z_6\) quotient & anomaly algebra; Witten anomaly; CKM CP counting & realized one-generation/one-Higgs package and MAR admissibility premises & Phase I realized-branch theorem/corollary chain \\
\end{longtable}
}
\endgroup

The recovered core in this appendix is therefore D1--D5 together with D7--D9. D6 appears only to keep the Phase-I / Phase-II boundary explicit. Nothing in this appendix should be read as promoting D10, D11, or D12 to theorem level.
\subsection{A.2 Secondary Quantitative and Continuation Benchmarks}\label{a2-quantitative-predictions}

This appendix does not extend the theorem-level list. It records every non-core branch that is still useful to track:

\begin{enumerate}
\item \textbf{Phase II input-dependent corollary (D6):} the capacity relation for \(\Lambda\) and the associated order-of-magnitude capacity / neutrino estimate.
\item \textbf{Phase II current implementation: calibration sector (D10):} gauge/electroweak closure of the present supplement-backed implementation, including the associated beta-shift comparison used in the current pipeline.
\item \textbf{Phase II current implementation: secondary quantitative branch (D11):} the Higgs/top critical-surface branch once the D10 trajectory and the extra critical-surface condition are imposed.
\item \textbf{Phase III deferred continuations (D12 and beyond):} charged-lepton/Koide, texture, neutrino mass/mixing constructions beyond the D6 estimate, dark-sector, baryogenesis, black-hole spectroscopy, string/worldsheet, proton-spin, and any proton-lifetime estimate beyond the structural exclusion of gauge-mediated decay.
\end{enumerate}

Failure of item 4 retracts the corresponding continuation only. Failure of item 3 challenges the D11 branch; failure of item 2 challenges the D10 implementation; failure of item 1 challenges the screen-capacity identification only. None of those failures by itself erases the recovered core D1--D5 plus D7--D9.

\subsection{A.3 Future Directions}\label{a3-future-directions}

\begin{enumerate}
\item construct or prove the scaling-limit regime required by the D3--D5 relativity chain, including the null-stress bridge and fixed-cap stationarity premises;
\item extend the compact-gauge reconstruction from the bosonic internal-gauge branch to the fermionic/super-Tannakian chiral branch while preserving the D7--D9 Standard Model structural chain;
\item derive or independently fix the external continuous inputs $P$ and $N_{\mathrm{scr}}$, and complete a fully auditable D10--D11 numerical pipeline with explicit running, threshold, and matching conventions;
\item separate the D6 capacity relation from flavor-style add-ons by deciding exactly which neutrino or cosmological continuations survive without additional ans\"atze;
\item only after items 1--4 are secured, reopen the Phase-III continuations: charged-lepton/Koide structure, texture exponents, dark-sector response laws, baryogenesis, black-hole spectroscopy, string/worldsheet reorganization, proton-spin bookkeeping, proton-lifetime estimates beyond the structural gauge-channel exclusion, and any similar late-stage continuation such as strong-CP proposals.
\end{enumerate}

Interpretive strange-loop / existence narratives sit below that ledger entirely: they are philosophical continuations, not appendix-backed theorem or benchmark claims.

Interpretive strange-loop / self-simulation narratives are not part of items 1--4 above. They are tracked, if at all, only as philosophical continuations rather than as quantitative supplement benchmarks. The supplement does not define a theorem-usable configuration space for a global emergence map, does not prove that observer formation and reconstruction close to an internal self-map on such a space, and does not supply compactness/contractivity/stability premises that would support an existence or stability theorem.

\begin{center}\rule{0.5\linewidth}{0.5pt}\end{center}

\section{Appendix B: Mixed-Status Numerical Benchmarks}\label{appendix-b-key-numerical-values}

{\def\LTcaptype{none} % do not increment counter
\begin{longtable}[]{@{}>{\RaggedRight\arraybackslash}p{0.240\textwidth}>{\RaggedRight\arraybackslash}p{0.240\textwidth}>{\RaggedRight\arraybackslash}p{0.240\textwidth}>{\RaggedRight\arraybackslash}p{0.240\textwidth}@{}}
\toprule\noalign{}
Quantity / branch & OPH Value & Observed & Status \\
\midrule\noalign{}
\endhead
\bottomrule\noalign{}
\endlastfoot
\(a_0\) (dark-sector continuation) & \(1.03 \times 10^{-10}\) m/s² & \(1.2 \times 10^{-10}\) m/s² & deferred continuation benchmark \\
\(\Delta\Sigma\) (spin continuation; QCD-dependent ansatz) & 0.308 & 0.29 ± 0.03 & deferred continuation benchmark \\
\(\eta_B\) (baryogenesis estimate) & \(4.6 \times 10^{-10}\) & \(6.1 \times 10^{-10}\) & deferred continuation benchmark \\
Koide \(Q\) (flavor continuation) & 2/3 & 0.666664 & deferred continuation benchmark \\
Koide \(\delta\) (flavor continuation) & 2/9 & 0.222225 & deferred continuation benchmark \\
\(\Delta b_3/\Delta b_2\) (D10 calibration) & 0.91 & 0.96 (MSSM) & supplement-backed calibration \\
\end{longtable}
}

\begin{center}\rule{0.5\linewidth}{0.5pt}\end{center}

\section{References}\label{references}

The derivations in this supplement are based on the OPH framework as specified in Part I of this manuscript. Standard results from:

\begin{itemize}
\tightlist
\item
  Tomita-Takesaki modular theory
\item
  Quantum Markov chain structure theorems (Hayden-Jozsa-Petz-Winter)
\item
  Fawzi-Renner recovery bounds
\item
  Gleason\textquotesingle s theorem
\item
  Jacobson\textquotesingle s entanglement equilibrium
\item
  Peter-Weyl decomposition
\item
  Heat kernel asymptotics on compact Lie groups
\item
  Tannaka-Krein reconstruction
\item
  \textquotesingle t Hooft anomaly matching and instanton calculus
\end{itemize}
